\subsection{Matched working risk helps conditionally, rather than universally}
The second development grid separates additional predictor flexibility from acquisition changes (Table~\ref{tab:step10matched}; Fig.~\ref{fig:step10task}). With the same expanded F1 kernel, five-label historical random-query MAE is 3.317 for random supports, 2.863 for composition V, and 2.792 for centered squared-risk selection. Finite-intercept GP selection gives 2.790, essentially the same result; removing its finite intercept prior therefore does not account for the point improvement. Centered absolute-risk selection gives 2.845 and does not improve on squared-risk selection. The expanded predictor is selected from block kernels in all twenty-five histories, but its improvement over the same-tuning isotropic predictor is small: under composition V, the two errors are 2.863 and 2.897. Kernel-family selection frequency alone is not a performance result.

The locked covariance-averaged absolute-risk candidate has historical MAE 3.030. None of its five conditional family-adjusted comparisons supports superiority; composition V and robust multi-view V give 2.863 and 2.859. Covariance averaging and a loss-matched objective do not supply a validated new-algorithm advantage. The primary family remains intact, while secondary contrasts and all predictor combinations are reported in the SI.

\begin{table}[htbp]\centering\small
\caption{Five-label F1 comparisons with predictor parameters held fixed across policies. MAE is in percentage points. History is development over 22 sources; both publication applications use the same previously opened family. The catalyst query and pool have disjoint reported catalyst identifiers.}\label{tab:step10matched}
\begin{tabular}{llrrr}\toprule
Acquisition & Predictor & History & Opened random & Catalyst holdout \\ \midrule
Random & Expanded kernel & 3.317 & 4.848 & 4.915 \\
Composition V & Expanded kernel & 2.863 & 4.164 & 5.055 \\
Pool V & Expanded kernel & 3.053 & 3.978 & 4.578 \\
Finite-intercept GP & Expanded kernel & 2.790 & 5.460 & 4.199 \\
Centered squared risk & Expanded kernel & 2.792 & 5.460 & 4.199 \\
Centered absolute risk & Expanded kernel & 2.845 & 5.788 & 4.199 \\
Averaged absolute risk & Expanded kernel & 3.030 & 4.793 & 4.381 \\
Random & Strong local RBF & 3.288 & 4.792 & 4.818 \\
Composition V & Strong local RBF & 2.851 & 3.813 & 4.680 \\
Centered squared risk & Strong local RBF & 2.784 & 5.145 & 3.689 \\
\bottomrule\end{tabular}
\end{table}

\plotfigure{step10_01_task_conditioned}{Selection and predictor effects under distinct query tasks at five labels. Within a task, each predictor's settings are fixed across all policies. The first panel is historical development; the other panels use the same opened publication. Better historical working-risk acquisition does not preserve its ranking under every query partition.}{fig:step10task}

The catalyst stress query comprises 435 measurements for 29 catalyst identifiers, with all 1304 support-pool rows belonging to the other 87 identifiers. In contrast, the earlier random partition has all 116 identities represented in both query and pool. Under catalyst holdout, centered squared-risk selection gives expanded-kernel MAE 4.199, compared with 4.915 for random and 5.055 for composition V. On exactly those same support sets, the strong isotropic RBF is better still: 3.689, compared with 4.818 for random and 4.680 for composition V. This is a 23.4\% point improvement over random for an unchanged established predictor. Finite-intercept, centered squared and centered absolute selection choose the same five supports in this partition; the benefit cannot be assigned specifically to free centering or absolute risk.

The regular random query gives the opposite ordering. Expanded-kernel centered squared selection has MAE 5.460, worse than random at 4.848 and composition V at 4.164; the corresponding strong-RBF errors are 5.145, 4.792 and 3.813. The wider kernel family also fails to improve on that RBF in the opened applications. Thus acquisition must be evaluated for the deployed predictor and query task; a lower working variance is not evidence of a universal reduction in actual error. These are single-family stress/development point estimates, without a population confidence interval. Reported-identifier separation does not certify unseen synthesis, physical batches or laboratories.

All 25 new historical comparisons and all three opened-publication partitions were executed. They add 82,350 and 1,830 evaluation cells, respectively, and 9,795,380 total scalar query predictions. Every scalar loss and all 16,836 support orders are audited; separate kernel/sklearn checks reconstruct 825,840 query values. Mathematical tests also independently verify the free-intercept weights and working-risk covariance. No global backbones were newly fitted in this extension: the earlier frozen models are reused, while kernel tuning, designs and predictions are new executed work.
